El proyecto SHIFTLY contempla el desarrollo de un prototipo de una
aplicación web orientada al aprendizaje del idioma inglés. La propuesta
busca ofrecer una experiencia de aprendizaje organizada y adaptable a los
intereses y necesidades del usuario. Como parte del alcance inicial, la aplicación contempla el registro e inicio de sesión del usuario, la realización de un diagnóstico para determinar su nivel de inglés y la selección de un área de aprendizaje.
Las áreas de aprendizaje consideradas son:
\begin{itemize}
\item Programación
\item Gastronomía
\item Turismo
\item Negocios
\item Inglés cotidiano
\end{itemize}
Después de seleccionar su área de interés, el usuario podrá acceder a
diferentes contenidos de aprendizaje, incluyendo lecciones, vocabulario y
ejercicios interactivos. También se contempla la evaluación de las
actividades y el seguimiento del progreso del usuario. El proyecto incluye diferentes interfaces que representan las principales funciones de la aplicación, entre ellas el diagnóstico inicial, la selección de áreas de aprendizaje, el dashboard, las tarjetas de
vocabulario, los ejercicios interactivos, el mapa de lecciones, la biblioteca y repaso, los resultados y las recomendaciones de aprendizaje.
El desarrollo corresponde a un prototipo inicial por lo tanto, algunas
funcionalidades se presentan principalmente a nivel visual y de
interacción de acuerdo con los objetivos establecidos para esta primer etapa del
proyecto